\section{Introduction}
\label{sec:introduction}

The automated analysis of digitised histopathology slides is a high-impact application of
computer vision. Pathologists must routinely identify and localise the tissue compartments
present in a biopsy: adipose tissue, mucus, stroma, lymphocytes, tumour epithelium, and
several others, a task that is time-consuming and subject to inter-observer variability. Deep
convolutional neural networks have demonstrated expert-level accuracy on such tasks
\citep{kather2019predicting}, but their data requirements remain a practical barrier: annotating
histopathology images at patch level requires specialist expertise and is therefore expensive.

A structural feature that distinguishes histopathology images from natural photographs is the
absence of a canonical orientation. A glass slide can be placed on the microscope stage at any
angle; the resulting whole-slide image therefore has no preferred ``up'' direction. This is not
a merely philosophical observation, it has a concrete consequence for learning. When a
standard CNN is trained on such data, its filters must implicitly learn to respond to the same
tissue pattern at every orientation, effectively duplicating parameter usage across rotations
\citep{veeling2018rotation}. Alternatively, the practitioner augments training data with
random rotations and reflections, forcing the network to encounter each orientation during
training, but providing no architectural guarantee of equivariance at test time.

Group-equivariant convolutional neural networks (\gcnn{}s), introduced by
\citet{cohen2016group}, offer a principled alternative. By replacing standard convolutions
with group convolutions defined over symmetry groups that include rotations and reflections,
the network is equivariant to those transformations \emph{by construction}. Each filter is
effectively shared across all orientations in the group, so the representational capacity
of the network is spent on learning content rather than orientation, leading to improved
parameter efficiency. In the digital pathology context this inductive bias is especially
natural: \citet{veeling2018rotation} showed that \gcnn{}s substantially improve tumour
detection on lymph-node metastases data.

This project investigates the following research question:
\begin{quote}
  \emph{Do group-equivariant CNNs provide a sample-efficiency or accuracy advantage over
  standard CNNs, with and without augmentation, on multi-class colorectal tissue
  classification, and how does this advantage scale with the quantity of labelled training
  data?}
\end{quote}

This question is addressed through a controlled experimental study on the NCT-CRC-HE-100K
dataset \citep{kather2018100}, comparing three model conditions across multiple training-data
fractions. The experimental pipeline is implemented in PyTorch~\citep{paszke2019pytorch}
and configured via Hydra~\citep{yadan2019hydra}, ensuring full reproducibility.

Our main contributions are:
\begin{itemize}
  \item A systematic comparison of a \gcnn{} (\(p4m\) group) against a standard CNN and an
    augmentation-trained CNN on a nine-class histopathology benchmark.
  \item A data-scaling study covering training-set fractions from 5\% to 100\%, designed to
    isolate where the equivariant inductive bias is most beneficial.
  \item An analysis of per-class behaviour, illuminating which tissue types benefit most from
    rotational equivariance.
\end{itemize}